\section{Compléments sur les faisceaux quasi-cohérents}
\label{section:I.9-fr}

\subsection{Produit tensoriel de faisceaux quasi-cohérents.}
\label{subsection:I.9.1-fr}

\begin{proposition}[9.1.1]
\label{I.9.1.1-fr}
Soit $X$ un préschéma (resp. un préschéma localement noethérien). Soient
$\mathscr{F}$ et $\mathscr{G}$ deux $\mathscr{O}_X$-Modules quasi-cohérents
(resp. cohérents)~; alors
$\mathscr{F}\otimes_{\mathscr{O}_X}\mathscr{G}$ est quasi-cohérent (resp.
cohérent) et de type fini si $\mathscr{F}$ et $\mathscr{G}$ sont de type
fini. Si $\mathscr{F}$ admet une présentation finie et si $\mathscr{G}$ est
quasi-cohérent (resp. cohérent), alors
$\mathscr{H}\!om_{\mathscr{O}_X}(\mathscr{F},\mathscr{G})$ est quasi-cohérent
(resp. cohérent).
\end{proposition}

La question étant locale, on peut supposer $X$ affine (resp. affine
noethérien)~; en outre, si $\mathscr{F}$ est cohérent, on peut supposer qu'il
est le conoyau d'un homomorphisme
$\mathscr{O}_X^m\to\mathscr{O}_X^n$. Les assertions relatives aux faisceaux
quasi-cohérents résultent alors de
(\hyperref[I.1.3.12-fr]{1.3.12}) et
(\hyperref[I.1.3.9-fr]{1.3.9})~; les assertions relatives aux faisceaux
cohérents résultent de (\hyperref[I.1.5.1-fr]{1.5.1}) et du fait que, si $M$
et $N$ sont des modules de type fini sur un anneau noethérien $A$,
$M\otimes_A N$ et $\operatorname{Hom}_A(M,N)$ sont des $A$-modules de type
fini.

\begin{definition}[9.1.2]
\label{I.9.1.2-fr}
Soient $X$, $Y$ deux $S$-préschémas, $p$, $q$ les projections de
$X\times_S Y$, $\mathscr{F}$ (resp. $\mathscr{G}$) un
$\mathscr{O}_X$-Module (resp. un $\mathscr{O}_Y$-Module) quasi-cohérent. On
appelle produit tensoriel de $\mathscr{F}$ et $\mathscr{G}$ sur
$\mathscr{O}_S$ (ou sur $S$) et on note
$\mathscr{F}\otimes_{\mathscr{O}_S}\mathscr{G}$ (ou
$\mathscr{F}\otimes_S\mathscr{G}$) le produit tensoriel
$p^*(\mathscr{F})\otimes_{\mathscr{O}_{X\times_S Y}}q^*(\mathscr{G})$ sur le
préschéma $X\times_S Y$.
\end{definition}

Si $X_i$ ($1\leq i\leq n$) sont des $S$-préschémas, $\mathscr{F}_i$ un
$\mathscr{O}_{X_i}$-Module quasi-cohérent ($1\leq i\leq n$), on définit de
la même manière le produit tensoriel
$\mathscr{F}_1\otimes_S\mathscr{F}_2\otimes_S\cdots\otimes_S\mathscr{F}_n$
sur le préschéma
$Z=X_1\times_S X_2\times_S\cdots\times_S X_n$~; c'est un
$\mathscr{O}_Z$-module \emph{quasi-cohérent} en vertu de
(\hyperref[I.9.1.1-fr]{9.1.1}) et de
(\hyperref[0.5.1.4-fr]{0, 5.1.4})~; il est \emph{cohérent} si les
$\mathscr{F}_i$ le sont et si $Z$ est \emph{localement noethérien} en vertu
de (\hyperref[I.9.1.1-fr]{9.1.1}),
(\hyperref[0.5.3.11-fr]{0, 5.3.11}) et
(\hyperref[I.6.1.1-fr]{6.1.1}).

On notera que si on prend $X=Y=S$, la définition
(\hyperref[I.9.1.2-fr]{9.1.2}) redonne le produit tensoriel de
$\mathscr{O}_S$-Modules. En outre, comme
$q^*(\mathscr{O}_Y)=\mathscr{O}_{X\times_S Y}$
(\hyperref[0.4.3.4-fr]{0, 4.3.4}), le produit
$\mathscr{F}\otimes_S\mathscr{O}_Y$ s'identifie canoniquement à
$p^*(\mathscr{F})$, et de même
$\mathscr{O}_X\otimes_S\mathscr{G}$ s'identifie canoniquement à
$q^*(\mathscr{G})$. Plus particulièrement, si on prend $Y=S$ et qu'on
désigne par $f$ le morphisme structural $X\to Y$, on a
$\mathscr{O}_X\otimes_Y\mathscr{G}=f^*(\mathscr{G})$~: le produit tensoriel
ordinaire et l'image réciproque apparaissent donc comme des cas particuliers
du produit tensoriel général.

La déf. (\hyperref[I.9.1.2-fr]{9.1.2}) entraîne immédiatement que, pour $X$
et $Y$ fixés, $\mathscr{F}\otimes_S\mathscr{G}$ est un
\emph{bifoncteur covariant additif et exact à droite} en $\mathscr{F}$ et
$\mathscr{G}$.

\begin{proposition}[9.1.3]
\label{I.9.1.3-fr}
Soient $S$, $X$, $Y$ trois schémas affines d'anneaux respectifs, $A$, $B$,
$C$, $B$ et $C$ étant des $A$-algèbres. Soit $M$ (resp. $N$) un $B$-module
(resp. $C$-module), $\mathscr{F}=\widetilde{M}$ (resp.
$\mathscr{G}=\widetilde{N}$) le faisceau quasi-cohérent associé~; alors
$\mathscr{F}\otimes_S\mathscr{G}$ est canoniquement isomorphe au faisceau
associé au $(B\otimes_A C)$-module $M\otimes_A N$.
\end{proposition}

\begin{proof}
\oldpage[I]{170}
En effet, en vertu de (\hyperref[I.1.6.5-fr]{1.6.5}),
$\mathscr{F}\otimes_S\mathscr{G}$ est canoniquement isomorphe au faisceau
associé au $(B\otimes_A C)$-module
\[
  \bigl(M\otimes_B(B\otimes_A C)\bigr)
  \otimes_{B\otimes_A C}
  \bigl((B\otimes_A C)\otimes_C N\bigr)
\]
et en raison des isomorphismes canoniques entre produits tensoriels, ce
dernier est isomorphe à
\[
  M\otimes_B(B\otimes_A C)\otimes_C N
  =(M\otimes_B B)\otimes_A(C\otimes_C N)=M\otimes_A N.
\]
\end{proof}

\begin{proposition}[9.1.4]
\label{I.9.1.4-fr}
Soient $f:T\to X$, $g:T\to Y$ deux $S$-morphismes, $\mathscr{F}$ (resp.
$\mathscr{G}$) un $\mathscr{O}_X$-Module (resp. $\mathscr{O}_Y$-Module)
quasi-cohérent. On a alors
\[
  (f,g)_S^*(\mathscr{F}\otimes_S\mathscr{G})
  =f^*(\mathscr{F})\otimes_{\mathscr{O}_T}g^*(\mathscr{G}).
\]
\end{proposition}

\begin{proof}
Si $p$, $q$ sont les projections de $X\times_S Y$, la formule résulte en
effet des relations $(f,g)_S^*\circ p^*=f^*$ et
$(f,g)_S^*\circ q^*=g^*$ (\hyperref[0.3.5.5-fr]{0, 3.5.5}), et du fait
qu'une image réciproque d'un produit tensoriel de faisceaux algébriques est
le produit tensoriel de leurs images réciproques
(\hyperref[0.4.3.3-fr]{0, 4.3.3}).
\end{proof}

\begin{corollary}[9.1.5]
\label{I.9.1.5-fr}
Soient $f:X\to X'$, $g:Y\to Y'$ deux $S$-morphismes, $\mathscr{F}'$ (resp.
$\mathscr{G}'$) un $\mathscr{O}_{X'}$-Module (resp.
$\mathscr{O}_{Y'}$-Module) quasi-cohérent. On a alors
\[
  (f\times_S g)^*(\mathscr{F}'\otimes_S\mathscr{G}')
  =f^*(\mathscr{F}')\otimes_S g^*(\mathscr{G}').
\]
\end{corollary}

\begin{proof}
Cela résulte de (\hyperref[I.9.1.4-fr]{9.1.4}) et du fait que
$f\times_S g=(f\circ p,g\circ q)_S$, $p$ et $q$ étant les projections de
$X\times_S Y$.
\end{proof}

\begin{corollary}[9.1.6]
\label{I.9.1.6-fr}
Soient $X$, $Y$, $Z$ trois $S$-préschémas, $\mathscr{F}$ (resp.
$\mathscr{G}$, $\mathscr{H}$) un $\mathscr{O}_X$-Module (resp.
$\mathscr{O}_Y$-Module, $\mathscr{O}_Z$-Module) quasi-cohérent~; le faisceau
$\mathscr{F}\otimes_S\mathscr{G}\otimes_S\mathscr{H}$ est l'image
réciproque de
$(\mathscr{F}\otimes_S\mathscr{G})\otimes_S\mathscr{H}$ par l'isomorphisme
canonique de $X\times_S Y\times_S Z$ sur
$(X\times_S Y)\times_S Z$.
\end{corollary}

\begin{proof}
En effet, cet isomorphisme s'écrit $(p_1,p_2)_S\times_S p_3$, en désignant
par $p_1$, $p_2$, $p_3$ les projections de $X\times_S Y\times_S Z$.

De même, l'image réciproque de $\mathscr{G}\otimes_S\mathscr{F}$ par
l'isomorphisme canonique de $X\times_S Y$ sur $Y\times_S X$ est
$\mathscr{F}\otimes_S\mathscr{G}$.
\end{proof}

\begin{corollary}[9.1.7]
\label{I.9.1.7-fr}
Si $X$ est un $S$-préschéma, tout $\mathscr{O}_X$-Module quasi-cohérent
$\mathscr{F}$ est l'image réciproque de
$\mathscr{F}\otimes_S\mathscr{O}_S$ par l'isomorphisme canonique de $X$ sur
$X\times_S S$ (\hyperref[I.3.3.3-fr]{3.3.3}).
\end{corollary}

\begin{proof}
En effet, cet isomorphisme est $(1_X,\varphi)_S$, où $\varphi$ est le
morphisme structural $X\to S$, et le corollaire résulte de
(\hyperref[I.9.1.4-fr]{9.1.4}) et du fait que
$\varphi^*(\mathscr{O}_S)=\mathscr{O}_X$.
\end{proof}

\begin{env}[9.1.8]
\label{I.9.1.8-fr}
Soient $X$ un $S$-préschéma, $\mathscr{F}$ un $\mathscr{O}_X$-Module
quasi-cohérent, $\varphi:S'\to S$ un morphisme~; on désigne par
$\mathscr{F}_{(\varphi)}$ ou $\mathscr{F}_{(S')}$ le faisceau
quasi-cohérent $\mathscr{F}\otimes_S\mathscr{O}_{S'}$, sur
$X\times_S S'=X_{(\varphi)}=X_{(S')}$~; donc
$\mathscr{F}_{(S')}=p^*(\mathscr{F})$, où $p$ est la projection
$X_{(S')}\to X$.
\end{env}

\begin{proposition}[9.1.9]
\label{I.9.1.9-fr}
Soit $\varphi':S''\to S'$ un morphisme. Pour tout
$\mathscr{O}_X$-Module quasi-cohérent $\mathscr{F}$ sur le $S$-préschéma
$X$, $(\mathscr{F}_{(\varphi)})_{(\varphi')}$ est l'image réciproque de
$\mathscr{F}_{(\varphi\circ\varphi')}$ par l'isomorphisme canonique
$(X_{(\varphi)})_{(\varphi')}\xrightarrow{\sim}
X_{(\varphi\circ\varphi')}$ (\hyperref[I.3.3.9-fr]{3.3.9}).
\end{proposition}

\begin{proof}
Cela résulte aussitôt des définitions et de
(\hyperref[I.3.3.9-fr]{3.3.9}), et s'écrit encore
\[
  (\mathscr{F}\otimes_S\mathscr{O}_{S'})
  \otimes_{S'}\mathscr{O}_{S''}
  =\mathscr{F}\otimes_S\mathscr{O}_{S''}.
  \tag{9.1.9.1}\label{I.9.1.9.1-fr}
\]
\end{proof}

\begin{proposition}[9.1.10]
\label{I.9.1.10-fr}
Soient $Y$ un $S$-préschéma, $f:X\to Y$ un $S$-morphisme. Pour tout
$\mathscr{O}_Y$-Module quasi-cohérent $\mathscr{G}$ et tout morphisme
$S'\to S$, on a
\[
  (f_{(S')})^*(\mathscr{G}_{(S')})
  =(f^*(\mathscr{G}))_{(S')}.
\]
\end{proposition}

\begin{proof}
Cela résulte aussitôt de la commutativité du diagramme
\oldpage[I]{171}
\[
  \xymatrix{
    X_{(S')}\ar[r]^{f_{(S')}}\ar[d] &
    Y_{(S')}\ar[d]\\
    X\ar[r]_f &
    Y
  }
\]
\end{proof}

\begin{corollary}[9.1.11]
\label{I.9.1.11-fr}
Soient $X$, $Y$ deux $S$-préschémas, $\mathscr{F}$ (resp. $\mathscr{G}$) un
$\mathscr{O}_X$-Module (resp. $\mathscr{O}_Y$-Module) quasi-cohérent.
L'image réciproque du faisceau
$(\mathscr{F}_{(S')}\otimes_{S'}\mathscr{G}_{(S')})$ par l'isomorphisme
canonique
$(X\times_S Y)_{(S')}\xrightarrow{\sim}
(X_{(S')})\times_{S'}(Y_{(S')})$
(\hyperref[I.3.3.10-fr]{3.3.10}) est égale à
$(\mathscr{F}\otimes_S\mathscr{G})_{(S')}$.
\end{corollary}

\begin{proof}
Si $p$, $q$ sont les projections de $X\times_S Y$, l'isomorphisme en
question n'est autre que $(p_{(S')},q_{(S')})_{S'}$~; le corollaire résulte
des prop. (\hyperref[I.9.1.4-fr]{9.1.4}) et
(\hyperref[I.9.1.10-fr]{9.1.10}).
\end{proof}

\begin{proposition}[9.1.12]
\label{I.9.1.12-fr}
Avec les notations de (\hyperref[I.9.1.2-fr]{9.1.2}) soient $z$ un point de
$X\times_S Y$, $x=p(z)$, $y=q(z)$~; la fibre
$(\mathscr{F}\otimes_S\mathscr{G})_z$ est isomorphe à
$(\mathscr{F}_x\otimes_{\mathscr{O}_x}\mathscr{O}_z)
\otimes_{\mathscr{O}_z}
(\mathscr{G}_y\otimes_{\mathscr{O}_y}\mathscr{O}_z)
=\mathscr{F}_x\otimes_{\mathscr{O}_x}\mathscr{O}_z
\otimes_{\mathscr{O}_y}\mathscr{G}_y$.
\end{proposition}

\begin{proof}
Comme on peut se ramener au cas affine, la proposition résulte de la formule
(\hyperref[I.1.6.5.1-fr]{1.6.5.1}).
\end{proof}

\begin{corollary}[9.1.13]
\label{I.9.1.13-fr}
Si $\mathscr{F}$ et $\mathscr{G}$ sont de type fini, on a
\[
  \operatorname{Supp}(\mathscr{F}\otimes_S\mathscr{G})
  =p^{-1}(\operatorname{Supp}(\mathscr{F}))
  \cap q^{-1}(\operatorname{Supp}(\mathscr{G})).
\]
\end{corollary}

\begin{proof}
Comme $p^*(\mathscr{F})$ et $q^*(\mathscr{G})$ sont de type fini sur
$\mathscr{O}_{X\times_S Y}$, on est ramené, par
(\hyperref[I.9.1.12-fr]{9.1.12}) et
(\hyperref[0.1.7.5-fr]{0, 1.7.5}), au cas particulier où
$\mathscr{G}=\mathscr{O}_Y$, c'est-à-dire à démontrer la formule
\[
  \operatorname{Supp}(p^{-1}(\mathscr{F}))
  =p^{-1}(\operatorname{Supp}(\mathscr{F})).
  \tag{9.1.13.1}\label{I.9.1.13.1-fr}
\]

Le même raisonnement que dans (\hyperref[0.1.7.5-fr]{0, 1.7.5}) ramène à
vérifier que l'on a, pour tout $z\in X\times_S Y$,
$\mathscr{O}_z/\mathfrak{m}_x\mathscr{O}_z\neq0$ (avec $x=p(z)$), ce qui
découle du fait que l'homomorphisme
$\mathscr{O}_x\to\mathscr{O}_z$ est \emph{local} par hypothèse.
\end{proof}

Nous laissons au lecteur le soin d'étendre à un produit d'un nombre
quelconque de facteurs les résultats démontrés dans ce numéro pour deux
facteurs.

\subsection{Image directe d'un faisceau quasi-cohérent.}
\label{subsection:I.9.2-fr}

\begin{proposition}[9.2.1]
\label{I.9.2.1-fr}
Soit $f:X\to Y$ un morphisme de préschémas. On suppose qu'il existe un
recouvrement $(Y_\alpha)$ de $Y$ par des ouverts affines ayant la propriété
suivante~: chacun des $f^{-1}(Y_\alpha)$ admet un recouvrement fini
$(X_{\alpha i})$ par des ouverts affines contenus dans
$f^{-1}(Y_\alpha)$, tel que chacune des intersections
$X_{\alpha i}\cap X_{\alpha j}$ soit elle-même réunion finie d'ouverts
affines. Dans ces conditions, pour tout $\mathscr{O}_X$-Module
quasi-cohérent $\mathscr{F}$, $f_*(\mathscr{F})$ est un
$\mathscr{O}_Y$-Module quasi-cohérent.
\end{proposition}

\begin{proof}
La question étant locale sur $Y$, on peut supposer $Y$ égal à l'un des
$Y_\alpha$, donc supprimer les indices $\alpha$.

\begin{enumerate}
\item[a)] Supposons d'abord que les $X_i\cap X_j$ soient eux-mêmes des
ouverts \emph{affines}. Posons
$\mathscr{F}_i=\mathscr{F}|X_i$,
$\mathscr{F}_{ij}=\mathscr{F}|(X_i\cap X_j)$, et soient
$\mathscr{F}'_i$ et $\mathscr{F}'_{ij}$ les images de $\mathscr{F}_i$ et
$\mathscr{F}_{ij}$ respectivement par les restrictions de $f$ à $X_i$ et
$X_i\cap X_j$~; on sait que les $\mathscr{F}'_i$ et
$\mathscr{F}'_{ij}$ sont quasi-cohérents
(\hyperref[I.1.6.3-fr]{1.6.3}). Posons
$\mathscr{G}=\bigoplus_i\mathscr{F}'_i$,
$\mathscr{H}=\bigoplus_{i,j}\mathscr{F}'_{ij}$~; $\mathscr{G}$ et
$\mathscr{H}$ sont des $\mathscr{O}_Y$-Modules quasi-cohérents~; nous allons
définir un homomorphisme $u:\mathscr{G}\to\mathscr{H}$ tel que
$f_*(\mathscr{F})$ soit le \emph{noyau} de $u$~; il en résultera que
$f_*(\mathscr{F})$ est quasi-cohérent
(\hyperref[I.1.3.9-fr]{1.3.9}). Il suffit de définir $u$
\oldpage[I]{172}
comme homomorphisme de préfaisceaux~; tenant compte des définitions de
$\mathscr{G}$ et $\mathscr{H}$, il suffit donc, pour tout ensemble ouvert
$W\subset Y$, de définir un homomorphisme
\[
  u_W:\bigoplus_i\Gamma(f^{-1}(W)\cap X_i,\mathscr{F})
  \longrightarrow
  \bigoplus_{i,j}\Gamma(f^{-1}(W)\cap X_i\cap X_j,\mathscr{F})
\]
de façon à satisfaire aux conditions de compatibilité usuelles lorsque $W$
varie. Si, pour toute section
$s_i\in\Gamma(f^{-1}(W)\cap X_i,\mathscr{F})$, on désigne par $s_{i|j}$ sa
restriction à $f^{-1}(W)\cap X_i\cap X_j$, on posera
\[
  u_W((s_i))=(s_{i|j}-s_{j|i})
\]
et les conditions de compatibilité sont remplies de façon évidente. Pour
prouver que le noyau $\mathscr{R}$ de $u$ est $f_*(\mathscr{F})$, définissons
un homomorphisme de $f_*(\mathscr{F})$ dans $\mathscr{R}$ en faisant
correspondre à toute section $s\in\Gamma(f^{-1}(W),\mathscr{F})$ la famille
$(s_i)$, où $s_i$ est la restriction de $s$ à $f^{-1}(W)\cap X_i$~; les
axiomes (F 1) et (F 2) des faisceaux (G, II, 1.1) entraînent que cet
homomorphisme est \emph{bijectif}, ce qui termine la démonstration dans ce cas.

\item[b)] Dans le cas général, le même raisonnement s'applique une fois que
l'on a établi que les $\mathscr{F}'_{ij}$ sont quasi-cohérents. Or, par
hypothèse, $X_i\cap X_j$ est réunion finie d'ouverts affines $X_{ijk}$~; et
comme les $X_{ijk}$ sont des ouverts affines \emph{dans un schéma},
l'intersection de deux quelconques d'entre eux est encore un ouvert affine
(\hyperref[I.5.5.6-fr]{5.5.6}). On est donc ramené au premier cas, et
(9.2.1) est donc démontrée.
\end{enumerate}
\end{proof}

\begin{corollary}[9.2.2]
\label{I.9.2.2-fr}
La conclusion de (9.2.1) est valable dans chacun des cas suivants~:
\begin{enumerate}
\item[a)] $f$ est séparé et quasi-compact.
\item[b)] $f$ est séparé et de type fini.
\item[c)] $f$ est quasi-compact et l'espace sous-jacent à $X$ est localement
noethérien.
\end{enumerate}
\end{corollary}

\begin{proof}
Dans le cas a), les $X_{\alpha i}\cap X_{\alpha j}$ sont affines
(\hyperref[I.5.5.6-fr]{5.5.6}). Le cas b) est un cas particulier de a)
(\hyperref[I.6.6.3-fr]{6.6.3}). Enfin, dans le cas c), on peut se ramener au
cas où $Y$ est affine et l'espace sous-jacent à $X$ noethérien~; alors $X$
admet un recouvrement ouvert affine fini $(X_i)$, et les $X_i\cap X_j$, étant
quasi-compacts, sont réunions finies d'ouverts affines
(\hyperref[I.2.1.3-fr]{2.1.3}).
\end{proof}

\subsection{Prolongement des sections de faisceaux quasi-cohérents.}
\label{subsection:I.9.3-fr}

\begin{theorem}[9.3.1]
\label{I.9.3.1-fr}
Soit $X$ un préschéma dont l'espace sous-jacent est noethérien, ou un schéma
dont l'espace sous-jacent est quasi-compact. Soient $\mathscr{L}$ un
$\mathscr{O}_X$-Module inversible
(\hyperref[0.5.4.1-fr]{0, 5.4.1}), $f$ une section de $\mathscr{L}$ au-dessus
de $X$, $X_f$ l'ensemble ouvert des $x\in X$ où $f(x)\neq0$
(\hyperref[0.5.5.1-fr]{0, 5.5.1}), $\mathscr{F}$ un
$\mathscr{O}_X$-Module quasi-cohérent.
\begin{enumerate}
\item[(i)] Si $s\in\Gamma(X,\mathscr{F})$ est telle que $s|X_f=0$, il existe
un entier $n>0$ tel que $s\otimes f^{\otimes n}=0$.
\item[(ii)] Pour toute section $s\in\Gamma(X_f,\mathscr{F})$, il existe un
entier $n>0$ tel que $s\otimes f^{\otimes n}$ se prolonge en une section de
$\mathscr{F}\otimes\mathscr{L}^{\otimes n}$ au-dessus de $X$.
\end{enumerate}
\end{theorem}

\begin{proof}
\begin{enumerate}
\item[(i)] Comme l'espace sous-jacent à $X$ est quasi-compact, donc réunion
finie d'ouverts affines $U_i$ tels que $\mathscr{L}|U_i$ soit isomorphe à
$\mathscr{O}_X|U_i$, on est ramené au cas où $X$ est affine et
$\mathscr{L}=\mathscr{O}_X$. Dans ce cas, $f$ s'identifie à un élément de
$A(X)$ et on a $X_f=D(f)$~; $s$ s'identifie à un élément d'un
$A(X)$-module $M$ et $s|X_f$ à l'élément correspondant de $M_f$, et le
résultat est trivial, compte tenu de la définition d'un module de fractions.
\oldpage[I]{173}

\item[(ii)] De nouveau $X$ est réunion finie d'ouverts affines $U_i$
($1\leq i\leq r$) tels que $\mathscr{L}|U_i\simeq\mathscr{O}_X|U_i$, et pour
chaque $i$,
$(s\otimes f^{\otimes n})|(U_i\cap X_f)$ s'identifie par l'isomorphisme
précédent à
$(f|(U_i\cap X_f))^n(s|(U_i\cap X_f))$. On sait alors
(\hyperref[I.1.4.1-fr]{1.4.1}) qu'il existe un entier $n>0$ tel que pour
chaque $i$, $(s\otimes f^{\otimes n})|(U_i\cap X_f)$ se prolonge en une
section $s_i$ de $\mathscr{F}\otimes\mathscr{L}^{\otimes n}$ au-dessus de
$U_i$. Soit $s_{i|j}$ la restriction de $s_i$ à $U_i\cap U_j$~; on a par
définition $s_{i|j}-s_{j|i}=0$ dans $X_f\cap U_i\cap U_j$. Or, si $X$ est un
espace noethérien, $U_i\cap U_j$ est quasi-compact~; si $X$ est un schéma,
$U_i\cap U_j$ est un ouvert affine
(\hyperref[I.5.5.6-fr]{5.5.6}), donc encore quasi-compact. En vertu de (i),
il existe donc un entier $m$ (indépendant de $i$ et $j$) tel que
$(s_{i|j}-s_{j|i})\otimes f^{\otimes m}=0$. On en conclut aussitôt qu'il
existe une section $s'$ de
$\mathscr{F}\otimes\mathscr{L}^{\otimes(n+m)}$ au-dessus de $X$, induisant
$s_i\otimes f^{\otimes m}$ au-dessus de chaque $U_i$, et induisant par suite
$s\otimes f^{\otimes(n+m)}$ au-dessus de $X_f$.
\end{enumerate}
\end{proof}

Les corollaires qui suivent donnent une interprétation du th. (9.3.1) en un
langage plus algébrique~:

\begin{corollary}[9.3.2]
\label{I.9.3.2-fr}
Les hypothèses étant celles de (9.3.1), considérons l'anneau gradué
$A_* = \Gamma_*(\mathscr{L})$ et le $A_*$-module gradué
$M_* = \Gamma_*(\mathscr{L},\mathscr{F})$
(\hyperref[0.5.4.6-fr]{0, 5.4.6}). Si $f\in A_n$, où $n\in\mathbb{Z}$, alors
on a un isomorphisme canonique
\[
  \Gamma(X_f,\mathscr{F})\xrightarrow{\sim}((M_*)_f)_0
\]
(sous-groupe du module de fractions $(M_*)_f$ formé des éléments de degré
$0$).
\end{corollary}

\begin{corollary}[9.3.3]
\label{I.9.3.3-fr}
On suppose vérifiées les hypothèses de (9.3.1), et on suppose en outre que
$\mathscr{L}=\mathscr{O}_X$. Alors si on pose
$A=\Gamma(X,\mathscr{O}_X)$, $M=\Gamma(X,\mathscr{F})$, le $A_f$-module
$\Gamma(X_f,\mathscr{F})$ est canoniquement isomorphe à $M_f$.
\end{corollary}

\begin{proposition}[9.3.4]
\label{I.9.3.4-fr}
Soient $X$ un préschéma noethérien, $\mathscr{F}$ un
$\mathscr{O}_X$-Module cohérent, $\mathscr{J}$ un faisceau cohérent d'idéaux
dans $\mathscr{O}_X$ tels que le support de $\mathscr{F}$ soit contenu dans
celui de $\mathscr{O}_X/\mathscr{J}$. Alors il existe un entier $n>0$ tel que
$\mathscr{J}^n\mathscr{F}=0$.
\end{proposition}

\begin{proof}
Comme $X$ est réunion finie d'ouverts affines dont les anneaux sont
noethériens, on peut supposer $X$ affine d'anneau $A$ noethérien~; alors
$\mathscr{F}=\widetilde M$, où $M=\Gamma(X,\mathscr{F})$ est un $A$-module
de type fini, et $\mathscr{J}=\widetilde{\mathfrak{J}}$, où
$\mathfrak{J}=\Gamma(X,\mathscr{J})$ est un idéal de $A$
(\hyperref[I.1.4.1-fr]{1.4.1} et \hyperref[I.1.5.1-fr]{1.5.1}). Comme $A$
est noethérien, $\mathfrak{J}$ admet un système fini de générateurs $f_i$
($1\leq i\leq m$). Par hypothèse, toute section de $\mathscr{F}$ au-dessus
de $X$ est nulle dans chacun des $D(f_i)$~; si $s_j$ ($1\leq j\leq q$) sont
des sections de $\mathscr{F}$ engendrant $M$, il existe donc un entier $h$
indépendant de $i$ et $j$, tel que $f_i^h s_j=0$
(\hyperref[I.1.4.1-fr]{1.4.1}), donc $f_i^h s=0$ pour tout $s\in M$. On en
conclut que si $n=mh$, on a $\mathfrak{J}^nM=0$, et par suite le
$\mathscr{O}_X$-Module correspondant
$\mathscr{J}^n\mathscr{F}=\widetilde{\mathfrak{J}^nM}$
(\hyperref[I.1.3.13-fr]{1.3.13}) est nul.
\end{proof}

\begin{corollary}[9.3.5]
\label{I.9.3.5-fr}
Sous les hypothèses de (9.3.4), il existe un sous-préschéma fermé $Y$ de
$X$, dont l'espace sous-jacent est le support de
$\mathscr{O}_X/\mathscr{J}$, et tel que, si $j:Y\to X$ est l'injection
canonique, on ait $\mathscr{F}=j_*(j^*(\mathscr{F}))$.
\end{corollary}

\begin{proof}
Notons d'abord que les supports de $\mathscr{O}_X/\mathscr{J}$ et de
$\mathscr{O}_X/\mathscr{J}^n$ sont les mêmes car si
$\mathscr{J}_x=\mathscr{O}_x$, on a aussi
$\mathscr{J}_x^n=\mathscr{O}_x$, et on a d'autre part
$\mathscr{J}_x^n\subset\mathscr{J}_x$ pour tout $x\in X$. On peut donc en
vertu de (9.3.4) supposer que $\mathscr{J}\mathscr{F}=0$~; on prend alors
pour $Y$ le sous-préschéma fermé de $X$ défini par $\mathscr{J}$, et comme
$\mathscr{F}$ est alors un $(\mathscr{O}_X/\mathscr{J})$-Module, la
conclusion est immédiate.
\end{proof}

\oldpage[I]{174}

\subsection{Prolongement des faisceaux quasi-cohérents.}
\label{subsection:I.9.4-fr}

\begin{env}[9.4.1]
\label{I.9.4.1-fr}
Soient $X$ un espace topologique, $\mathscr{F}$ un faisceau d'ensembles
(resp. de groupes, d'anneaux) sur $X$, $U$ une partie ouverte de $X$,
$\psi:U\to X$ l'injection canonique, $\mathscr{G}$ un sous-faisceau de
$\mathscr{F}|U=\psi^*(\mathscr{F})$. Comme $\psi_*$ est exact à gauche,
$\psi_*(\mathscr{G})$ est un sous-faisceau de
$\psi_*(\psi^*(\mathscr{F}))$~; si on considère l'homomorphisme canonique
$\rho:\mathscr{F}\to\psi_*(\psi^*(\mathscr{F}))$
(\hyperref[0.3.5.3-fr]{0, 3.5.3}), nous désignerons par
$\overline{\mathscr{G}}$ le sous-faisceau
$\rho^{-1}(\psi_*(\mathscr{G}))$ de $\mathscr{F}$. Il résulte immédiatement
des définitions que pour tout ouvert $V$ de $X$,
$\Gamma(V,\overline{\mathscr{G}})$ est formé des sections
$s\in\Gamma(V,\mathscr{F})$ dont la restriction à $V\cap U$ soit une section
de $\mathscr{G}$ au-dessus de $V\cap U$. On a donc
$\overline{\mathscr{G}}|U=\psi^*(\overline{\mathscr{G}})=\mathscr{G}$, et
$\overline{\mathscr{G}}$ est le \emph{plus grand} sous-faisceau de
$\mathscr{F}$ induisant $\mathscr{G}$ sur $U$~; nous dirons que
$\overline{\mathscr{G}}$ est le \emph{prolongement canonique} du
sous-faisceau $\mathscr{G}$ de $\mathscr{F}|U$ en un sous-faisceau de
$\mathscr{F}$.
\end{env}

\begin{proposition}[9.4.2]
\label{I.9.4.2-fr}
Soient $X$ un préschéma, $U$ une partie ouverte de $X$ telle que l'injection
canonique $j:U\to X$ soit un morphisme quasi-compact
\emph{(ce qui sera vérifié pour \emph{tout} $U$ si l'espace sous-jacent $X$
est \emph{localement noethérien}
(\hyperref[I.6.6.4-fr]{6.6.4, (i)}))}. Alors~:
\begin{enumerate}
\item[(i)] Pour tout $(\mathscr{O}_X|U)$-Module quasi-cohérent
  $\mathscr{G}$, $j_*(\mathscr{G})$ est un $\mathscr{O}_X$-Module
  quasi-cohérent et on a
  $j_*(\mathscr{G})|U=j^*(j_*(\mathscr{G}))=\mathscr{G}$.
\item[(ii)] Pour tout $\mathscr{O}_X$-Module quasi-cohérent $\mathscr{F}$ et
  tout sous-$(\mathscr{O}_X|U)$-Module quasi-cohérent $\mathscr{G}$, le
  prolongement canonique $\overline{\mathscr{G}}$ de $\mathscr{G}$ (9.4.1)
  est un sous-$\mathscr{O}_X$-Module quasi-cohérent de $\mathscr{F}$.
\end{enumerate}
\end{proposition}

\begin{proof}
Si $j=(\psi,\theta)$ ($\psi$ étant l'injection $U\to X$ des espaces
sous-jacents), on a par définition $j_*(\mathscr{G})=\psi_*(\mathscr{G})$
pour tout $(\mathscr{O}_X|U)$-Module $\mathscr{G}$, et en outre
$j^*(\mathscr{H})=\psi^*(\mathscr{H})=\mathscr{H}|U$ pour tout
$\mathscr{O}_X$-Module $\mathscr{H}$ en raison de la définition d'un
préschéma induit sur un ouvert. (i) est donc un cas particulier de
(\hyperref[I.9.2.2-fr]{9.2.2, a)})~; pour la même raison,
$j_*(j^*(\mathscr{F}))$ est quasi-cohérent, et comme
$\overline{\mathscr{G}}$ est l'image réciproque de $j_*(\mathscr{G})$ par
l'homomorphisme $\rho:\mathscr{F}\to j_*(j^*(\mathscr{F}))$, (ii) résulte de
(\hyperref[I.4.1.1-fr]{4.1.1}).
\end{proof}

On notera que l'hypothèse que le morphisme $j:U\to X$ est quasi-compact est
aussi vérifiée lorsque l'ouvert $U$ est \emph{quasi-compact} et $X$ un
\emph{schéma}~: en effet, $U$ est alors réunion finie d'ouverts affines
$U_i$, et pour tout ouvert affine $V$ de $X$, $V\cap U_i$ est un ouvert
affine (\hyperref[I.5.5.6-fr]{5.5.6}), donc quasi-compact.

\begin{corollary}[9.4.3]
\label{I.9.4.3-fr}
Soient $X$ un préschéma, $U$ un ouvert quasi-compact de $X$ tel que le
morphisme d'injection $j:U\to X$ soit quasi-compact. Supposons en outre que
tout $\mathscr{O}_X$-Module quasi-cohérent soit limite inductive de ses
sous-$\mathscr{O}_X$-Modules quasi-cohérents de type fini
\emph{(ce qui a lieu lorsque $X$ est un \emph{schéma affine})}. Soient alors
$\mathscr{F}$ un $\mathscr{O}_X$-Module quasi-cohérent, $\mathscr{G}$ un
sous-$(\mathscr{O}_X|U)$-Module quasi-cohérent de type fini de
$\mathscr{F}|U$. Il existe alors un sous-$\mathscr{O}_X$-Module
quasi-cohérent de type fini $\mathscr{G}'$ de $\mathscr{F}$ tel que
$\mathscr{G}'|U=\mathscr{G}$.
\end{corollary}

\begin{proof}
En effet, on a $\mathscr{G}=\overline{\mathscr{G}}|U$, et
$\overline{\mathscr{G}}$ est quasi-cohérent d'après (9.4.2), donc limite
inductive de ses sous-$\mathscr{O}_X$-Modules quasi-cohérents de type fini
$\mathscr{H}_\lambda$. Par suite $\mathscr{G}$ est limite inductive des
$\mathscr{H}_\lambda|U$, donc égal à un des $\mathscr{H}_\lambda|U$
puisqu'il est de type fini (\hyperref[0.5.2.3-fr]{0, 5.2.3}).
\end{proof}

\begin{remark}[9.4.4]
\label{I.9.4.4-fr}
Supposons que pour \emph{tout} ouvert affine $U\subset X$ le morphisme
d'injection $U\to X$ soit quasi-compact. Alors si la conclusion de (9.4.3) a
lieu pour tout ouvert affine $U$ et tout sous-$(\mathscr{O}_X|U)$-Module
quasi-cohérent de type fini $\mathscr{G}$ de $\mathscr{F}|U$,
\oldpage[I]{175}
il en résulte que $\mathscr{F}$ est limite inductive de ses
sous-$\mathscr{O}_X$-Modules quasi-cohérents de type fini. En effet, pour
tout ouvert affine $U\subset X$, on a $\mathscr{F}|U=\widetilde{M}$, où $M$
est un $A(U)$-module, et comme ce dernier est limite inductive de ses
sous-modules de type fini, $\mathscr{F}|U$ est limite inductive de ses
sous-$(\mathscr{O}_X|U)$-Modules quasi-cohérents de type fini
(\hyperref[I.1.3.9-fr]{1.3.9}). Or, par hypothèse, chacun de ces sous-Modules
est induit sur $U$ par un sous-$\mathscr{O}_X$-Module quasi-cohérent de type
fini $\mathscr{G}_{\lambda,U}$ de $\mathscr{F}$. Les sommes finies des
$\mathscr{G}_{\lambda,U}$ sont encore des $\mathscr{O}_X$-Modules
quasi-cohérents de type fini, car la question est locale et le cas où $X$
est affine a été traité dans (\hyperref[I.1.3.10-fr]{1.3.10})~; il est clair
alors que $\mathscr{F}$ est limite inductive de ces sommes finies, d'où notre
assertion.
\end{remark}

\begin{corollary}[9.4.5]
\label{I.9.4.5-fr}
Sous les hypothèses de (9.4.3), pour tout $(\mathscr{O}_X|U)$-Module
quasi-cohérent de type fini $\mathscr{G}$, il existe un
$\mathscr{O}_X$-Module quasi-cohérent de type fini $\mathscr{G}'$ tel que
$\mathscr{G}'|U=\mathscr{G}$.
\end{corollary}

\begin{proof}
Comme $\mathscr{F}=j_*(\mathscr{G})$ est quasi-cohérent (9.4.2) et
$\mathscr{F}|U=\mathscr{G}$, il suffit d'appliquer (9.4.3) à $\mathscr{F}$.
\end{proof}

\begin{lemma}[9.4.6]
\label{I.9.4.6-fr}
Soient $X$ un préschéma, $L$ un ensemble bien ordonné,
$(V_\lambda)_{\lambda\in L}$ un recouvrement de $X$ par des ouverts affines,
$U$ un ouvert de $X$~; pour tout $\lambda\in L$, on pose
$W_\lambda=\bigcup_{\mu<\lambda}V_\mu$. On suppose que~:
\begin{enumerate}
  \item[1\textsuperscript{o}] Pour tout $\lambda\in L$,
    $V_\lambda\cap W_\lambda$ soit quasi-compact~;
  \item[2\textsuperscript{o}] Le morphisme d'immersion $U\to X$ est
    quasi-compact.
\end{enumerate}
Alors, pour tout $\mathscr{O}_X$-Module quasi-cohérent $\mathscr{F}$ et tout
sous-$(\mathscr{O}_X|U)$-Module quasi-cohérent de type fini $\mathscr{G}$ de
$\mathscr{F}|U$, il existe un sous-$\mathscr{O}_X$-Module quasi-cohérent de
type fini $\mathscr{G}'$ de $\mathscr{F}$ tel que
$\mathscr{G}'|U=\mathscr{G}$.
\end{lemma}

\begin{proof}
Posons $U_\lambda=U\cup W_\lambda$~; on va définir par récurrence transfinie
une famille $(\mathscr{G}'_\lambda)$, où $\mathscr{G}'_\lambda$ est un
sous-$(\mathscr{O}_X|U_\lambda)$-Module quasi-cohérent de type fini de
$\mathscr{F}|U_\lambda$, tel que
$\mathscr{G}'_\lambda|U_\mu=\mathscr{G}'_\mu$ pour $\mu<\lambda$ et
$\mathscr{G}'_\lambda|U=\mathscr{G}$. L'unique sous-$\mathscr{O}_X$-Module
$\mathscr{G}'$ de $\mathscr{F}$ tel que
$\mathscr{G}'|U_\lambda=\mathscr{G}'_\lambda$ pour tout $\lambda\in L$
(\hyperref[0.3.3.1-fr]{0, 3.3.1}) répondra à la question. Supposons donc les
$\mathscr{G}'_\mu$ définis et ayant les propriétés précédentes pour
$\mu<\lambda$~; si $\lambda$ n'a pas de prédécesseur on prendra pour
$\mathscr{G}'_\lambda$ l'unique sous-$(\mathscr{O}_X|U_\lambda)$-Module de
$\mathscr{F}|U_\lambda$ tel que
$\mathscr{G}'_\lambda|U_\mu=\mathscr{G}'_\mu$ pour tout $\mu<\lambda$, ce
qui est licite puisque les $U_\mu$ avec $\mu<\lambda$ forment alors un
recouvrement de $U_\lambda$. Si au contraire $\lambda=\mu+1$, on a
$U_\lambda=U_\mu\cup V_\mu$, et il suffira de définir un
sous-$(\mathscr{O}_X|V_\mu)$-Module quasi-cohérent de type fini
$\mathscr{G}''_\mu$ de $\mathscr{F}|V_\mu$ tel que
\[
  \mathscr{G}''_\mu|(U_\mu\cap V_\mu)
  =\mathscr{G}'_\mu|(U_\mu\cap V_\mu)~;
\]
on prendra ensuite pour $\mathscr{G}'_\lambda$ le
sous-$(\mathscr{O}_X|U_\lambda)$-Module de $\mathscr{F}|U_\lambda$ tel que
$\mathscr{G}'_\lambda|U_\mu=\mathscr{G}'_\mu$ et
$\mathscr{G}'_\lambda|V_\mu=\mathscr{G}''_\mu$
(\hyperref[0.3.3.1-fr]{0, 3.3.1}). Or, comme $V_\mu$ est affine, l'existence
de $\mathscr{G}''_\mu$ est assurée par (9.4.3) dès que l'on a démontré que
$U_\mu\cap V_\mu$ est quasi-compact~; mais $U_\mu\cap V_\mu$ est réunion de
$U\cap V_\mu$ et de $W_\mu\cap V_\mu$, qui sont tous deux quasi-compacts en
vertu de l'hypothèse.
\end{proof}

\begin{theorem}[9.4.7]
\label{I.9.4.7-fr}
Soient $X$ un préschéma, $U$ un ouvert de $X$. On suppose vérifiée l'une des
conditions suivantes~:
\begin{enumerate}
  \item[(a)] L'espace sous-jacent à $X$ est localement noethérien.
  \item[(b)] $X$ est un schéma quasi-compact et $U$ un ouvert quasi-compact.
\end{enumerate}
Pour tout $\mathscr{O}_X$-Module quasi-cohérent $\mathscr{F}$ et tout
sous-$(\mathscr{O}_X|U)$-Module quasi-cohérent de type fini $\mathscr{G}$ de
$\mathscr{F}|U$, il existe alors un sous-$\mathscr{O}_X$-Module
quasi-cohérent de type fini $\mathscr{G}'$ de $\mathscr{F}$ tel que
$\mathscr{G}'|U=\mathscr{G}$.
\end{theorem}

\oldpage[I]{176}

\begin{proof}
Soit $(V_\lambda)_{\lambda\in L}$ un recouvrement de $X$ par des ouverts
affines, $L$ étant supposé fini dans le cas b)~; $L$ étant muni d'une
structure d'ensemble bien ordonné, il suffit de vérifier que les conditions
du lemme (9.4.6) sont satisfaites. C'est évident dans l'hypothèse a), les
espaces $V_\lambda$ étant noethériens. Dans l'hypothèse b), les
$V_\lambda\cap V_\mu$ sont affines (\hyperref[I.5.5.6-fr]{5.5.6}), donc
quasi-compacts, et comme $L$ est fini, $V_\lambda\cap W_\lambda$ est
quasi-compact. D'où le théorème.
\end{proof}

\begin{corollary}[9.4.8]
\label{I.9.4.8-fr}
Sous les hypothèses de (9.4.7), pour tout $(\mathscr{O}_X|U)$-Module
quasi-cohérent de type fini $\mathscr{G}$, il existe un
$\mathscr{O}_X$-Module quasi-cohérent de type fini $\mathscr{G}'$ tel que
$\mathscr{G}'|U=\mathscr{G}$.
\end{corollary}

\begin{proof}
Il suffit d'appliquer (9.4.7) à $\mathscr{F}=j_*(\mathscr{G})$, qui est
quasi-cohérent (9.4.2) et tel que $\mathscr{F}|U=\mathscr{G}$.
\end{proof}

\begin{corollary}[9.4.9]
\label{I.9.4.9-fr}
Soit $X$ un préschéma dont l'espace sous-jacent est localement noethérien, ou
un schéma quasi-compact. Alors tout $\mathscr{O}_X$-Module quasi-cohérent est
limite inductive de ses sous-$\mathscr{O}_X$-Modules quasi-cohérents de type
fini.
\end{corollary}

\begin{proof}
Cela résulte de (9.4.7) et de la remarque (9.4.4).
\end{proof}

\begin{corollary}[9.4.10]
\label{I.9.4.10-fr}
Sous les hypothèses de (9.4.9), si un $\mathscr{O}_X$-Module quasi-cohérent
$\mathscr{F}$ est tel que tout sous-$\mathscr{O}_X$-Module quasi-cohérent de
type fini de $\mathscr{F}$ soit engendré par ses sections au-dessus de $X$,
alors $\mathscr{F}$ est engendré par ses sections au-dessus de $X$.
\end{corollary}

\begin{proof}
En effet, soient $U$ un voisinage ouvert affine d'un point $x\in X$, et soit
$s$ une section de $\mathscr{F}$ au-dessus de $U$~; le
sous-$\mathscr{O}_X$-Module $\mathscr{G}$ de $\mathscr{F}|U$ engendré par
$s$ est quasi-cohérent et de type fini, donc il existe un
sous-$\mathscr{O}_X$-Module quasi-cohérent de type fini $\mathscr{G}'$ de
$\mathscr{F}$ tel que $\mathscr{G}'|U=\mathscr{G}$ (9.4.7). Par hypothèse,
il y a donc un nombre fini de sections $t_i$ de $\mathscr{G}'$ au-dessus de
$X$ et des sections $a_i$ de $\mathscr{O}_X$ au-dessus d'un voisinage
$V\subset U$ de $x$ telles que
$s|V=\sum_i a_i\mathbin{\cdot}(t_i|V)$, ce qui démontre le corollaire.
\end{proof}

\subsection{Image fermée d'un préschéma. Adhérence d'un sous-préschéma.}
\label{subsection:I.9.5-fr}

\begin{proposition}[9.5.1]
\label{I.9.5.1-fr}
Soit $f:X\to Y$ un morphisme de préschémas tel que $f_*(\mathscr{O}_X)$ soit
un $\mathscr{O}_Y$-Module quasi-cohérent (ce qui a lieu si $f$ est
quasi-compact et si de plus $f$ est séparé ou $X$ localement noethérien
(\hyperref[I.9.2.2-fr]{9.2.2})). Alors il existe un plus petit sous-préschéma
$Y'$ de $Y$ tel que l'injection canonique $j:Y'\to Y$ majore $f$ (ou, ce qui
revient au même (\hyperref[I.4.4.1-fr]{4.4.1}), tel que le sous-préschéma
$f^{-1}(Y')$ de $X$ soit \emph{identique} à $X$).
\end{proposition}

De façon plus précise~:

\begin{corollary}[9.5.2]
\label{I.9.5.2-fr}
Sous les conditions de (9.5.1), soit $f=(\psi,\theta)$ et soit $\mathscr{J}$
le noyau (quasi-cohérent) de l'homomorphisme
$\theta:\mathscr{O}_Y\to f_*(\mathscr{O}_X)$. Alors le sous-préschéma fermé
$Y'$ de $Y$ défini par $\mathscr{J}$ vérifie les conditions de (9.5.1).
\end{corollary}

\begin{proof}
Comme le foncteur $\psi^*$ est exact, la factorisation canonique
$\theta:\mathscr{O}_Y\to\mathscr{O}_Y/\mathscr{J}
\xrightarrow{\theta'}\psi_*(\mathscr{O}_X)$ donne
(\hyperref[0.3.5.4.3-fr]{0, 3.5.4.3}) une factorisation
\[
  \theta^\sharp:\psi^*(\mathscr{O}_Y)
  \to\psi^*(\mathscr{O}_Y)/\psi^*(\mathscr{J})
  \xrightarrow{{\theta'}^\sharp}\mathscr{O}_X~;
\]
comme pour tout $x\in X$, $\theta_x^\sharp$ est un homomorphisme local, il
en est de même de ${\theta'_x}^\sharp$~; si on désigne par $\psi_0$
l'application continue $\psi$ considérée comme application de $X$ dans
$X'$, par $\theta_0$ la restriction
$\theta'|X':(\mathscr{O}_Y/\mathscr{J})|X'
\to\psi_*(\mathscr{O}_X)|X'=(\psi_0)_*(\mathscr{O}_X)$, on voit donc que
$f_0=(\psi_0,\theta_0)$ est un morphisme de préschémas $X\to X'$
(\hyperref[I.2.2.1-fr]{2.2.1}) tel que $f=j\circ f_0$. Si maintenant $X''$
\oldpage[I]{177}
est un second sous-préschéma fermé de $Y$, défini par un faisceau
quasi-cohérent d'idéaux $\mathscr{J}'$ de $\mathscr{O}_Y$, et tel que
l'injection $j':X''\to Y$ majore $f$, on doit tout d'abord avoir
$X''\supset\psi(X)$, donc $X'\subset X''$ puisque $X''$ est fermé. En outre,
pour tout $y\in X''$, $\theta$ doit se factoriser en
$\mathscr{O}_y\to\mathscr{O}_y/\mathscr{J}'_y
\to(\psi_*(\mathscr{O}_X))_y$, ce qui par définition entraîne
$\mathscr{J}'_y\subset\mathscr{J}_y$, et par suite $X'$ est un
sous-préschéma fermé de $X''$ (\hyperref[I.4.1.10-fr]{4.1.10}).
\end{proof}

\begin{definition}[9.5.3]
\label{I.9.5.3-fr}
Lorsqu'il existe un plus petit sous-préschéma fermé $Y'$ de $Y$ tel que
l'injection canonique $j:Y'\to Y$ majore $f$, on dit que $Y'$ est le
préschéma image fermée de $X$ par le morphisme $f$.
\end{definition}

\begin{proposition}[9.5.4]
\label{I.9.5.4-fr}
Si $f_*(\mathscr{O}_X)$ est un $\mathscr{O}_Y$-Module quasi-cohérent,
l'espace sous-jacent de l'image fermée de $X$ par $f$ est l'adhérence
$\overline{f(X)}$ dans $Y$.
\end{proposition}

\begin{proof}
Comme le support de $f_*(\mathscr{O}_X)$ est contenu dans
$\overline{f(X)}$, on a (avec les notations de
(\hyperref[I.9.5.2-fr]{9.5.2})) $\mathscr{J}_y=\mathscr{O}_y$ pour
$y\notin\overline{f(X)}$, donc le support de
$\mathscr{O}_Y/\mathscr{J}$ est contenu dans $\overline{f(X)}$. En outre, ce
support est fermé et contient $f(X)$~: en effet, si $y\in f(X)$, l'élément
unité de l'anneau $(\psi_*(\mathscr{O}_X))_y$ n'est pas nul, étant le germe
en $y$ de la section
\[
  1\in\Gamma(X,\mathscr{O}_X)=\Gamma(Y,\psi_*(\mathscr{O}_X))~;
\]
comme il est l'image par $\theta$ de l'élément unité de $\mathscr{O}_y$, ce
dernier n'appartient pas à $\mathscr{J}_y$, donc
$\mathscr{O}_y/\mathscr{J}_y\neq0$~; ceci achève la démonstration.
\end{proof}

\begin{proposition}[9.5.5]
\label{I.9.5.5-fr}
\emph{(transitivité des images fermées).} Soient $f:X\to Y$ et $g:Y\to Z$
deux morphismes de préschémas~; on suppose que l'image fermée $Y'$ de $X$
par $f$ existe, et que, si $g'$ est la restriction de $g$ à $Y'$, l'image
fermée $Z'$ de $Y'$ par $g'$ existe. Alors l'image fermée de $X$ par
$g\circ f$ existe et est égale à $Z'$.
\end{proposition}

\begin{proof}
Il suffit (\hyperref[I.9.5.1-fr]{9.5.1}) de montrer que $Z'$ est le plus
petit sous-préschéma fermé $Z_1$ de $Z$ tel que le sous-préschéma fermé
$(g\circ f)^{-1}(Z_1)$ de $X$ (égal à $f^{-1}(g^{-1}(Z_1))$ par
(\hyperref[I.4.4.2-fr]{4.4.2})) soit égal à $X$~; il revient au même de dire
que $Z'$ est le plus petit sous-préschéma fermé de $Z$ tel que $f$ soit
majoré par l'injection $g^{-1}(Z_1)\to Y$
(\hyperref[I.4.4.1-fr]{4.4.1}). Or, en vertu de l'existence de l'image
fermée $Y'$, tout $Z_1$ ayant cette propriété est tel que $g^{-1}(Z_1)$
majore $Y'$, ce qui équivaut à dire que
$j^{-1}(g^{-1}(Z_1))=g'^{-1}(Z_1)=Y'$, en désignant par $j$ l'injection
$Y'\to Y$. Par définition de $Z'$, on en conclut bien que $Z'$ est le plus
petit sous-préschéma fermé de $Z$ vérifiant la condition précédente.
\end{proof}

\begin{corollary}[9.5.6]
\label{I.9.5.6-fr}
Soit $f:X\to Y$ un $S$-morphisme tel que $Y$ soit l'image fermée de $X$ par
$f$. Soit $Z$ un $S$-schéma~; si deux $S$-morphismes $g_1$, $g_2$ de $Y$
dans $Z$ sont tels que $g_1\circ f=g_2\circ f$, alors $g_1=g_2$.
\end{corollary}

\begin{proof}
Soit $h=(g_1,g_2)_S:Y\to Z\times_S Z$~; comme la diagonale
$T=\Delta_Z(Z)$ est un sous-préschéma fermé de $Z\times_S Z$,
$Y'=h^{-1}(T)$ est un sous-préschéma fermé de $Y$
(\hyperref[I.4.4.1-fr]{4.4.1}). Posons $u=g_1\circ f=g_2\circ f$~; on a
alors par définition du produit $h'=h\circ f=(u,u)_S$, donc
$h\circ f=\Delta_Z\circ u$~; comme $\Delta_Z^{-1}(T)=Z$, on a
$h'^{-1}(T)=u^{-1}(Z)=X$, donc $f^{-1}(Y')=X$. On en conclut
(\hyperref[I.4.4.1-fr]{4.4.1}) que l'injection canonique $Y'\to Y$ majore
$f$, donc $Y'=Y$ par hypothèse~; par suite
(\hyperref[I.4.4.1-fr]{4.4.1}), $h$ se factorise en $\Delta_Z\circ v$, où
$v$ est un morphisme $Y\to Z$, ce qui entraîne $g_1=g_2=v$.
\end{proof}

\begin{remark}[9.5.7]
\label{I.9.5.7-fr}
Si $X$ et $Y$ sont des $S$-schémas, la prop.
(\hyperref[I.9.5.6-fr]{9.5.6}) signifie que
\oldpage[I]{178}
lorsque $Y$ est l'image fermée de $X$ par $f$, $f$ est un
\emph{épimorphisme} dans la catégorie des $S$-schémas (T, 1.1). Nous
démontrerons au chap. V que, réciproquement, si l'image fermée $Y'$ de $X$
par $f$ existe et si $f$ est un épimorphisme de $S$-schémas, alors on a
nécessairement $Y'=Y$.
\end{remark}

\begin{proposition}[9.5.8]
\label{I.9.5.8-fr}
Supposons vérifiées les hypothèses de (\hyperref[I.9.5.1-fr]{9.5.1}), et
soit $Y'$ l'image fermée de $X$ par $f$. Pour tout ouvert $V$ de $Y$, soit
$f_V:f^{-1}(V)\to V$ la restriction de $f$~; alors l'image fermée de
$f^{-1}(V)$ par $f_V$ dans $V$ existe et est égale au préschéma induit par
$Y'$ sur l'ouvert $V\cap Y'$ de $Y'$ (autrement dit, au sous-préschéma
$\inf(V,Y')$ de $Y$ (\hyperref[I.4.4.3-fr]{4.4.3})).
\end{proposition}

\begin{proof}
Posons $X'=f^{-1}(V)$~; comme l'image directe de $\mathscr{O}_{X'}$ par
$f_V$ n'est autre que la restriction de $f_*(\mathscr{O}_X)$ à $V$, il est
clair que le noyau $\mathscr{J}'$ de l'homomorphisme
$\mathscr{O}_V\to(f_V)_*(\mathscr{O}_{X'})$ est la restriction de
$\mathscr{J}$ à $V$, d'où aussitôt la proposition.
\end{proof}

On notera que ce résultat s'interprète en disant que la formation de l'image
fermée commute à une extension $Y_1\to Y$ du préschéma de base, qui est une
\emph{immersion ouverte}. Nous verrons au chap. IV qu'il en est de même pour
une extension $Y_1\to Y$ qui est un morphisme \emph{plat}, pourvu que $f$
soit séparé et quasi-compact.

\begin{proposition}[9.5.9]
\label{I.9.5.9-fr}
Soit $f:X\to Y$ un morphisme tel que l'image fermée $Y'$ de $X$ par $f$
existe.
\begin{enumerate}
  \item[(i)] Si $X$ est réduit, il en est de même de $Y'$.
  \item[(ii)] Si on suppose vérifiées les hypothèses de
    (\hyperref[I.9.5.1-fr]{9.5.1}) et si $X$ est irréductible (resp.
    intègre), il en est de même de $Y'$.
\end{enumerate}
\end{proposition}

\begin{proof}
Par hypothèse, le morphisme $f$ se factorise en
$X\xrightarrow{g}Y'\xrightarrow{j}Y$, où $j$ est l'injection canonique.
Comme $X$ est réduit, $g$ se factorise en
$X\xrightarrow{h}Y'_{\mathrm{red}}\xrightarrow{j'}Y'$, où $j'$ est
l'injection canonique (\hyperref[I.5.2.2-fr]{5.2.2}), et il résulte alors de
la définition de $Y'$ que $Y'_{\mathrm{red}}=Y'$. Si de plus les conditions
de (\hyperref[I.9.5.1-fr]{9.5.1}) sont remplies, il résulte de
(\hyperref[I.9.5.4-fr]{9.5.4}) que $f(X)$ est dense dans $Y'$~; si $X$ est
irréductible, il en est donc de même de $Y'$
(\hyperref[0.2.1.5-fr]{0, 2.1.5}). L'assertion relative aux préschémas
intègres résulte de la conjonction des deux autres.
\end{proof}

\begin{proposition}[9.5.10]
\label{I.9.5.10-fr}
Soit $Y$ un sous-préschéma d'un préschéma $X$, tel que l'injection canonique
$i:Y\to X$ soit un morphisme quasi-compact. Il existe alors un plus petit
sous-préschéma fermé $\overline{Y}$ de $X$ majorant $Y$~; son espace
sous-jacent est l'adhérence de celui de $Y$~; ce dernier est ouvert dans son
adhérence, et le préschéma $Y$ est induit sur cet ouvert par $\overline{Y}$.
\end{proposition}

\begin{proof}
Il suffit d'appliquer (\hyperref[I.9.5.1-fr]{9.5.1}) à l'injection $j$, qui
est séparée (\hyperref[I.5.5.1-fr]{5.5.1}) et quasi-compacte par hypothèse~;
(\hyperref[I.9.5.1-fr]{9.5.1}) prouve donc l'existence de $\overline{Y}$ et
(\hyperref[I.9.5.4-fr]{9.5.4}) montre que son espace sous-jacent est
l'adhérence de $Y$ dans $X$~; comme $Y$ est localement fermé dans $X$, il
est ouvert dans $\overline{Y}$, et la dernière assertion provient de
(\hyperref[I.9.5.8-fr]{9.5.8}) appliqué à un ouvert $V$ de $X$ tel que $Y$
soit fermé dans $V$.
\end{proof}

Avec ces notations, si l'injection $V\to X$ est quasi-compacte, et si
$\mathscr{J}$ est le faisceau quasi-cohérent d'idéaux de
$\mathscr{O}_X|V$ définissant le sous-préschéma fermé $Y$ de $V$, il résulte
de (\hyperref[I.9.5.1-fr]{9.5.1}) que le faisceau quasi-cohérent d'idéaux de
$\mathscr{O}_X$ définissant $\overline{Y}$ est le prolongement canonique
(\hyperref[I.9.4.1-fr]{9.4.1}) $\overline{\mathscr{J}}$ de $\mathscr{J}$,
car c'est évidemment le plus grand sous-faisceau d'idéaux quasi-cohérent de
$\mathscr{O}_X$ induisant $\mathscr{J}$ sur $V$.
\oldpage[I]{179}
\begin{corollary}[9.5.11]
\label{I.9.5.11-fr}
Sous les hypothèses de (\hyperref[I.9.5.10-fr]{9.5.10}), toute section de
$\mathscr{O}_{\overline{Y}}$ au-dessus d'un ouvert $V$ de $\overline{Y}$
qui est nulle dans $V\cap Y$ est nulle.
\end{corollary}

\begin{proof}
En vertu de (\hyperref[I.9.5.8-fr]{9.5.8}), on peut se ramener au cas où
$V=\overline{Y}$. Si l'on tient compte de ce que les sections de
$\mathscr{O}_{\overline{Y}}$ au-dessus de $\overline{Y}$ correspondent
canoniquement aux $\overline{Y}$-sections de
$\overline{Y}\otimes_{\mathbb{Z}}\mathbb{Z}[T]$
(\hyperref[I.3.3.15-fr]{3.3.15}) et de ce que ce dernier est séparé sur
$\overline{Y}$, le corollaire apparaît comme un cas particulier de
(\hyperref[I.9.5.6-fr]{9.5.6}).
\end{proof}

Lorsqu'il existe un plus petit sous-préschéma fermé $\overline{Y}$ de $X$
majorant un sous-préschéma $Y$ de $X$, on dit que $\overline{Y}$ est
l'\emph{adhérence} de $Y$ dans $X$, lorsqu'il n'en résulte pas de
confusion.

\subsection{Faisceaux quasi-cohérents d'algèbres~; changement du faisceau
structural.}
\label{subsection:I.9.6-fr}

\begin{proposition}[9.6.1]
\label{I.9.6.1-fr}
Soient $X$ un préschéma, $\mathscr{B}$ une $\mathscr{O}_X$-Algèbre
quasi-cohérente (\hyperref[0.5.1.3-fr]{0, 5.1.3}). Pour qu'un
$\mathscr{B}$-Module $\mathscr{F}$ soit quasi-cohérent (sur l'espace
annelé $(X,\mathscr{B})$), il faut et il suffit que $\mathscr{F}$ soit un
$\mathscr{O}_X$-Module quasi-cohérent.
\end{proposition}

\begin{proof}
Comme la question est locale, on peut supposer $X$ affine d'anneau $A$, et
alors $\mathscr{B}=\widetilde{B}$, où $B$ est une $A$-algèbre
(\hyperref[I.1.4.3-fr]{1.4.3}). Si $\mathscr{F}$ est quasi-cohérent sur
l'espace annelé $(X,\mathscr{B})$, on peut aussi supposer que
$\mathscr{F}$ est le conoyau d'un $\mathscr{B}$-homomorphisme
$\mathscr{B}^{(I)}\to\mathscr{B}^{(J)}$~; comme cet homomorphisme est
aussi un $\mathscr{O}_X$-homomorphisme de $\mathscr{O}_X$-Modules, et que
$\mathscr{B}^{(I)}$ et $\mathscr{B}^{(J)}$ sont des
$\mathscr{O}_X$-Modules quasi-cohérents
(\hyperref[I.1.3.9-fr]{1.3.9, (ii)}), $\mathscr{F}$ est aussi un
$\mathscr{O}_X$-Module quasi-cohérent
(\hyperref[I.1.3.9-fr]{1.3.9, (i)}).

Inversement, si $\mathscr{F}$ est un $\mathscr{O}_X$-Module
quasi-cohérent, on a $\mathscr{F}=\widetilde{M}$, où $M$ est un $B$-module
(\hyperref[I.1.4.3-fr]{1.4.3})~; $M$ est isomorphe à un conoyau de
$B$-homomorphisme $B^{(I)}\to B^{(J)}$, donc $\mathscr{F}$ est un
$\mathscr{B}$-Module isomorphe au conoyau de l'homomorphisme correspondant
$\mathscr{B}^{(I)}\to\mathscr{B}^{(J)}$
(\hyperref[I.1.3.13-fr]{1.3.13}), ce qui achève la démonstration.
\end{proof}

En particulier, si $\mathscr{F}$ et $\mathscr{G}$ sont deux
$\mathscr{B}$-Modules quasi-cohérents,
$\mathscr{F}\otimes_{\mathscr{B}}\mathscr{G}$ est un
$\mathscr{B}$-Module quasi-cohérent~; il en est de même de
$\mathscr{H}\!om_{\mathscr{B}}(\mathscr{F},\mathscr{G})$ lorsqu'on
suppose en outre que $\mathscr{F}$ admet une présentation finie
(\hyperref[I.1.3.13-fr]{1.3.13}).

\begin{env}[9.6.2]
\label{I.9.6.2-fr}
Étant donné un préschéma $X$, nous dirons qu'une
$\mathscr{O}_X$-Algèbre quasi-cohérente $\mathscr{B}$ est de
\emph{type fini} si pour tout $x\in X$, il existe un voisinage ouvert
\emph{affine} $U$ de $x$ tel que $\Gamma(U,\mathscr{B})=B$ soit une
algèbre de type fini sur $\Gamma(U,\mathscr{O}_X)=A$. On a alors
$\mathscr{B}|U=\widetilde{B}$ et pour tout $f\in A$, la
$(\mathscr{O}_X|D(f))$-Algèbre induite $\mathscr{B}|D(f)$ est de type
fini, car elle est isomorphe à $\widetilde{B_f}$, et
$B_f=B\otimes_A A_f$ est évidemment une algèbre de type fini sur $A_f$.
Comme les $D(f)$ forment une base de la topologie de $U$, on en conclut que
si $\mathscr{B}$ est une $\mathscr{O}_X$-Algèbre quasi-cohérente de type
fini, pour tout ouvert $V$ de $X$, $\mathscr{B}|V$ est une
$(\mathscr{O}_X|V)$-Algèbre quasi-cohérente de type fini.
\end{env}

\begin{proposition}[9.6.3]
\label{I.9.6.3-fr}
Soit $X$ un préschéma localement noethérien. Toute
$\mathscr{O}_X$-Algèbre quasi-cohérente de type fini $\mathscr{B}$ est
alors un faisceau cohérent d'anneaux
(\hyperref[0.5.3.7-fr]{0, 5.3.7}).
\end{proposition}

\begin{proof}
On peut de nouveau se limiter au cas où $X$ est un schéma affine d'anneau
noethérien $A$, et où $\mathscr{B}=\widetilde{B}$, $B$ étant une
$A$-algèbre de type fini~; $B$ est donc un anneau noethérien. Cela étant, il
faut prouver que le noyau $\mathscr{N}$ d'un
$\mathscr{B}$-homomorphisme $\mathscr{B}^m\to\mathscr{B}$ est un
$\mathscr{B}$-
\oldpage[I]{180}
Module de type fini~; or, il est isomorphe (en tant que
$\mathscr{B}$-Module) à $\widetilde{N}$, où $N$ est le noyau de
l'homomorphisme correspondant de $B$-modules $B^m\to B$
(\hyperref[I.1.3.13-fr]{1.3.13}). Comme $B$ est noethérien, le sous-module
$N$ de $B^m$ est un $B$-module de type fini, donc il existe un homomorphisme
$B^p\to B^m$ d'image $N$~; la suite $B^p\to B^m\to B$ étant exacte, il
en est de même de la suite
$\mathscr{B}^p\to\mathscr{B}^m\to\mathscr{B}$ correspondante
(\hyperref[I.1.3.5-fr]{1.3.5}) et comme $\mathscr{N}$ est l'image de
$\mathscr{B}^p\to\mathscr{B}^m$
(\hyperref[I.1.3.9-fr]{1.3.9, (i)}), la proposition est démontrée.
\end{proof}

\begin{corollary}[9.6.4]
\label{I.9.6.4-fr}
Sous les hypothèses de (\hyperref[I.9.6.3-fr]{9.6.3}), pour qu'un
$\mathscr{B}$-Module $\mathscr{F}$ soit cohérent, il faut et il suffit
qu'il soit un $\mathscr{O}_X$-Module quasi-cohérent et un
$\mathscr{B}$-Module de type fini. S'il en est ainsi, et si $\mathscr{G}$
est un sous-$\mathscr{B}$-Module ou un $\mathscr{B}$-Module quotient de
$\mathscr{F}$, pour que $\mathscr{G}$ soit un $\mathscr{B}$-Module
cohérent, il faut et il suffit que $\mathscr{G}$ soit un
$\mathscr{O}_X$-Module quasi-cohérent.
\end{corollary}

\begin{proof}
Compte tenu de (\hyperref[I.9.6.1-fr]{9.6.1}), les conditions sur
$\mathscr{F}$ sont évidemment nécessaires~; montrons qu'elles sont
suffisantes. On peut se limiter au cas où $X$ est affine d'anneau
noethérien $A$, $\mathscr{B}=\widetilde{B}$, où $B$ est une $A$-algèbre de
type fini, $\mathscr{F}=\widetilde{M}$, où $M$ est un $B$-module, et où il
existe un $\mathscr{B}$-homomorphisme surjectif
$\mathscr{B}^m\to\mathscr{F}\to 0$. On a alors une suite exacte
correspondante $B^m\to M\to 0$, donc $M$ est un $B$-module de type fini~;
en outre, le noyau $P$ de l'homomorphisme $B^m\to M$ est alors un
$B$-module de type fini, puisque $B$ est noethérien. On en conclut
(\hyperref[I.1.3.13-fr]{1.3.13}) que $\mathscr{F}$ est le conoyau d'un
$\mathscr{B}$-homomorphisme $\mathscr{B}^n\to\mathscr{B}^m$, et est donc
cohérent puisque $\mathscr{B}$ est un faisceau cohérent d'anneaux
(\hyperref[0.5.3.4-fr]{0, 5.3.4}). Le même raisonnement montre qu'un
sous-$\mathscr{B}$-Module (resp. un $\mathscr{B}$-Module quotient)
quasi-cohérent de $\mathscr{F}$ est de type fini, d'où la seconde partie du
corollaire.
\end{proof}

\begin{proposition}[9.6.5]
\label{I.9.6.5-fr}
Soit $X$ un schéma quasi-compact ou un préschéma dont l'espace sous-jacent
est noethérien. Pour toute $\mathscr{O}_X$-Algèbre quasi-cohérente de type
fini $\mathscr{B}$, il existe un sous-$\mathscr{O}_X$-Module quasi-cohérent
de type fini $\mathscr{E}$ de $\mathscr{B}$ tel que $\mathscr{E}$ engendre
(\hyperref[0.4.1.4-fr]{0, 4.1.4}) la $\mathscr{O}_X$-Algèbre
$\mathscr{B}$.
\end{proposition}

\begin{proof}
En effet, il existe par hypothèse un recouvrement fini $(U_i)$ de $X$ formé
d'ouverts affines tels que $\Gamma(U_i,\mathscr{B})=B_i$ soit une algèbre
de type fini sur $\Gamma(U_i,\mathscr{O}_X)=A_i$~; soit $E_i$ un
sous-$A_i$-module de type fini de $B_i$ engendrant la $A_i$-algèbre $B_i$~;
en vertu de (\hyperref[I.9.4.7-fr]{9.4.7}), il existe un
sous-$\mathscr{O}_X$-Module $\mathscr{E}_i$ de $\mathscr{B}$,
quasi-cohérent et de type fini, tel que
$\mathscr{E}_i|U_i=\widetilde{E_i}$. Il est clair que la somme
$\mathscr{E}$ des $\mathscr{E}_i$ répond à la question.
\end{proof}

\begin{proposition}[9.6.6]
\label{I.9.6.6-fr}
Soit $X$ un préschéma dont l'espace sous-jacent est localement noethérien,
ou un schéma quasi-compact. Alors toute $\mathscr{O}_X$-Algèbre
quasi-cohérente $\mathscr{B}$ est limite inductive de ses
sous-$\mathscr{O}_X$-Algèbres quasi-cohérentes de type fini.
\end{proposition}

\begin{proof}
En effet, il résulte de (\hyperref[I.9.4.9-fr]{9.4.9}) que $\mathscr{B}$
est limite inductive (en tant que $\mathscr{O}_X$-Module) de ses
sous-$\mathscr{O}_X$-Modules quasi-cohérents de type fini~; ces derniers
engendrent des sous-$\mathscr{O}_X$-Algèbres quasi-cohérentes de type fini
de $\mathscr{B}$ (\hyperref[I.1.3.14-fr]{1.3.14}), dont $\mathscr{B}$ est
\emph{a fortiori} limite inductive.
\end{proof}
